
\section{Density Ratio Estimation and Reweighting}
\label{sec:dre}

% ---------------------------------------------------------------
% Sub-topics to cover:
%
% 1. The core statistical trick
%    - Given two distributions p(x) and q(x), we often want
%      the ratio r(x) = p(x)/q(x) without knowing either
%      density explicitly.
%    - Key insight: a binary classifier trained to distinguish
%      samples from p vs. q learns r(x) implicitly, because
%      the optimal classifier output is r(x)/(1 + r(x)).
%
% 2. Classifier-based density ratio estimation
%    - Train f_\theta(x) -> [0,1] with BCE loss.
%    - Recover r(x) = f(x)/(1 - f(x)).
%    - Works in high dimensions where direct density estimation fails.
%
% 3. Applications in HEP
%    - Monte Carlo reweighting: correct simulation to data.
%    - Importance sampling: reweight one generator to another.
%    - Omnifold (Ch.5): iterative density ratio estimation for
%      unbinned unfolding — this is the key application.
%    - CARL (Cranmer et al.) as a foundational reference.
%
% 4. Practical considerations
%    - Calibration of the classifier output is crucial.
%    - Extreme weights can cause variance blow-up: use weight
%      clipping or regularization.
%    - Choice of architecture matters: must have good coverage of
%      the support of both distributions.
